\documentclass[../../syllabus_1.tex]{subfiles}
\begin{document}
\raggedbottom
\chapter{Figures planes}

\section*{Matières}
\begin{enumerate}[cols=2, label=\textbf{\arabic*.}]
    \item Les types d'angles.
    \item Les triangles.
    \item Les quadrilatères.
    \item Droites remarquables.
\end{enumerate}

\newpage
\section*{Mots-clés}
\begin{enumerate}[cols=2, label=\textbf{\arabic*.}]
    \item Angles
    \item Rentrant
    \item Plein
    \item Adjacents
    \item Complémentaires
    \item Supplémentaires
    \item Opposés par le sommet
    \item Inégalité triangulaire
    \item Médiane
    \item Médiatrice
    \item Bissectrice
\end{enumerate}

\section*{Attendus}
\begin{enumerate}[label=\textbf{\arabic*.}]
   \item Identifier des figures simples (triangle, quadrilatère) sur base de caractéristiques : les côtés (nombre, longueur, parallélisme, perpendicularité) ; les angles (amplitude).
    \item Identifier des polygones réguliers sur base des caractéristiques : les côtés (nombre, longueur) ; les angles (amplitude).
    \item Associer aux angles particuliers (droit, aigu, obtus, plat, plein) leurs amplitudes.
    \item Identifier la médiatrice d’un segment, la bissectrice d’un angle.
    \item Identifier les droites remarquables d’un triangle.
    \item Définir la médiatrice d’un segment, la bissectrice d’un angle.
    \item Définir les droites remarquables d’un triangle.
    \item Identifier les quadrilatères admettant un centre de symétrie.
    \item Énoncer les propriétés des droites remarquables dans le triangle.
    \item Associer un symbole à l’objet qu’il désigne : point, segment, demi-droite, droite, angle, droites parallèles, droites perpendiculaires.
    \item Lire, interpréter et utiliser les notations, les symboles et le codage géométriques.
    \item Représenter une situation géométrique décrite à l’aide des notations, des symboles et du codage géométrique.
    \item Construire la médiatrice d’un segment, la bissectrice d’un angle.
    \item Construire les droites remarquables d’un triangle ou d’un quadrilatère.
    \item Mesurer l’amplitude d’un angle à l’aide du rapporteur.
    \item Construire un angle dont l’amplitude est donnée.
    \item Construire un triangle connaissant :\\
    - la longueur d’un côté et l’amplitude de ses deux angles adjacents ;\\
    - la longueur de deux côtés et l’amplitude de l’angle compris entre eux ;
    - la longueur des trois côtés.
    \item Construire un triangle isocèle ou équilatéral :\\
    - à la latte et au compas ;\\
    - à la latte et au rapporteur.
    \item Construire un losange connaissant la longueur du côté et l’amplitude d’un angle.
    \item Construire un parallélogramme connaissant la longueur de deux côtés et l’amplitude de l’angle compris entre eux.
    \item Construire une figure complexe (figure composée de figures simples), les étapes de construction étant données.
    \item Rédiger des étapes de construction d’une figure complexe (figure composée de figures simples) donnée.
    \item Identifier des angles complémentaires, supplémentaires, opposés par le sommet.
    \item Énoncer la propriété de la somme des amplitudes des angles intérieurs d’un triangle et d’un quadrilatère.
    \item Énoncer les propriétés relatives aux angles, aux axes et centre de symétrie, aux droites remarquables, dans des figures simples.
    \item Déterminer l’amplitude d’un angle en utilisant les propriétés des angles complémentaires, supplémentaires et opposés par le sommet.
    \item Déterminer l’amplitude d’un angle en utilisant la propriété relative à la somme des amplitudes des angles intérieurs d’un triangle et d’un quadrilatère.
    \item Terminer la construction d’une figure simple respectant des contraintes sur les droites remarquables.
    \item Terminer la construction d’une figure simple respectant des contraintes sur les axes et le centre de symétrie.
    \item Construire le cercle inscrit et le cercle circonscrit à un triangle.
    \item Justifier l’amplitude d’un angle en utilisant des propriétés.
    \item Construire une figure simple dont le centre de symétrie, le(les) axe(s) de symétrie sont donnés.
    \item Construire un quadrilatère (carré, rectangle, losange, parallélogramme) dont une diagonale ou une médiane est donnée.
    \item Construire un triangle dont deux médiatrices et un sommet sont donnés.
    \item Construire un triangle dont deux bissectrices, un sommet (appartenant à l’une d’elles) et l’amplitude de l’angle en ce sommet sont donnés.
    \item Construire un triangle isocèle dont une médiatrice (ou une bissectrice) et un sommet ou l’amplitude d’un angle ou la longueur d’un côté sont donnés.
    \item Construire un triangle équilatéral dont une médiatrice (ou une bissectrice) et un sommet ou la longueur d’un côté sont donnés.
    \item Résoudre un problème mobilisant des propriétés relatives aux angles et justifier.
    \item Exprimer le périmètre d’un triangle équilatéral, d’un rectangle, d’un carré, d’un cercle en fonction de ses dimensions.
    \item Justifier les liens entre les formules d’aire du rectangle, du parallélogramme, du trapèze, du triangle et du losange.
    \item Énoncer la formule du calcul de l’aire du disque.
    \item Exprimer l’aire d’une figure simple (rectangle, parallélogramme, trapèze, triangle, losange, carré) après avoir repéré les dimensions utiles.
    \item Calculer la longueur d’un côté d’un polygone régulier (triangle équilatéral, carré, pentagone, hexagone, octogone, décagone) dont le périmètre est donné.
    \item Calculer le périmètre d’un cercle.
    \item Calculer l’aire d’une figure simple (triangle, quadrilatère, disque).
    \item Décomposer une surface complexe en plusieurs surfaces connues (triangles, quadrilatères, disques) pour en calculer l’aire, avec des formules préalablement construites.
    \item Résoudre des problèmes faisant intervenir des calculs d’aire et de périmètre de figures simples ou complexes, en situations contextualisées.
    \item Représenter une figure simple dont l’aire est identique à celle d’une autre figure simple donnée. Exemple : Représenter un rectangle dont l’aire est identique à celle d’un triangle donné.
\end{enumerate}

\end{document}